% \vspace*{-0.3\baselineskip}
\section{Background and Related Works}

% Graph Chain-of-Thought (Graph-CoT) is a reasoning framework proposed in~\cite{jin2024graph} that augments LLMs with graph-based context to support iterative reasoning. Each iteration in the Graph-CoT framework involves three sub-steps: LLM reasoning, interaction with the graph, and graph-based execution. Through this iterative process, the LLM performs \textbf{chain-of-thought} reasoning while selectively retrieving information from the \textbf{Graph RAG}. The loop continues until the LLM reaches a final conclusion during the reasoning step.

\subsection{Problem Definition}

With the proliferation of large-scale graph data across diverse real-world applications, an increasing number of studies leverage graphs to enhance reasoning capabilities in LLMs \cite{wu2024can, cao2024graphreason, jin2024graph}.
A graph reasoning framework enables LLMs to iteratively reason over graph-structured data, allowing them to access domain-specific knowledge and improve answer accuracy.
% while maintaining efficiency in token consumption and response latency.

\label{sec:problem-definition}
\stitle{Graph Reasoning Problem.} 
Let $q \sim \mathcal{D}$ denote a query drawn from a distribution $\mathcal{D}$, and let $G$ be an external graph that can be accessed during reasoning. 
We consider a set of agents $\mathcal{A}=\{a_1,\dots,a_K\}$ coordinated by a policy $\pi$. 
Given a query $q$, the interaction among the policy $\pi$, the agents $\mathcal{A}$, and the graph $G$ induces a stochastic reasoning trajectory
$\tau \sim P^{\pi,\mathcal{A}}(\,\cdot \mid q, G)$.
An answer $\hat{y}$ is generated from the trajectory through a decoding function $g$, i.e., $\hat{y}=g(\tau)$.
Correctness is evaluated by an accuracy function $Acc(\hat{y}, y^\star(q)) \in [0,1]$, where $y^\star(q)$ denotes the ground-truth answer.
Efficiency is measured by the total token cost $Tok(\tau)$ and end-to-end latency $Lat(\tau)$.

The optimization objective is to design a policy $\pi$ and agent set $\mathcal{A}$ that maximize expected accuracy under bounded cost and latency constraints:
\begin{equation}
\begin{array}{cl}
\max_{\pi,\,\mathcal{A}} 
& \mathbb{E}_{q \sim \mathcal{D}}\, \mathbb{E}_{\tau \sim P^{\pi,\mathcal{A}}(\,\cdot \mid q, G)} \big[ Acc(g(\tau), y^{\star}(q)) \big] \\[3pt]
\text{s.t.} 
& \mathbb{E}_{q, \tau}[Tok(\tau)] \leq \tau_{\max}, \\[2pt]
& \mathbb{E}_{q, \tau}[Lat(\tau)] \leq \ell_{\max}.
\end{array}
\end{equation}

In practice, a trajectory $\tau=(s_0, u_0, o_1, \dots, s_T)$ records the sequence of states $s_t$, actions $u_t$ (e.g., selecting an agent or issuing a graph operation), and observations $o_{t+1}$ returned by either an agent $a_k \in \mathcal{A}$ or the graph $G$. 
The policy $\pi$ governs routing decisions, agent selection, memory management, and termination.
The total token cost and latency are instantiated as:
$
Tok(\tau)=\sum_{t=0}^{T}\!\big(\mathrm{tok\_in}(u_t)+\mathrm{tok\_out}(o_{t+1})\big), 
Lat(\tau)=\sum_{t=0}^{T}\!\Delta t_t.
$

\noindent\textit{Remark.} Unlike most prior work, which pays limited attention to token cost and efficiency (i.e., end-to-end latency and throughput), we explicitly model both via $\mathrm{Tok}(\tau)$ and $\mathrm{Lat}(\tau)$ and treat them as top-priority optimization objectives in our design.



\subsection{Existing Graph Reasoning Methods} 

Recent research has explored how to design effective policies and agent architectures to improve reasoning accuracy over graph-structured data~\cite{li2024graphreader, wang2023knowledgpt, jiang2023structgpt, jin2024graph}. 
\textsc{GraphReader}~\cite{li2024graphreader} adopts a planning-based, coarse-to-fine exploration strategy that leverages predefined functions to read node contents and neighbourhoods, while maintaining iterative note-taking and reflection. 
\textsc{KnowledGPT}~\cite{wang2023knowledgpt} formulates graphs and queries as executable programs, using a Program-of-Thought framework to generate and execute search code that interacts with a lightweight knowledge-base API for multi-hop reasoning. 
\textsc{StructGPT}~\cite{jiang2023structgpt} employs an Iterative Reading–then–Reasoning (IRR) loop, in which task-specific interfaces extract graph evidence that the LLM subsequently reasons over after linearization. 
However, these approaches primarily emphasize reasoning accuracy, often overlooking efficiency metrics such as token consumption and end-to-end latency.

Among existing methods, \textsc{Graph-CoT}~\cite{jin2024graph} represents the current state-of-the-art solution for graph reasoning. 
It instantiates an iterative \emph{Reasoning–Interaction–Execution} framework, enabling the LLM to autonomously determine subsequent graph operations, issue primitive function calls, execute them on an external graph, and repeat this process until convergence. 
This explicit stepwise traversal allows the model to exploit graph topology dynamically, contrasting with static retrieval-augmented generation (RAG) approaches that rely on context stuffing without structural awareness.

\subsection{RAG and CoT}

\stitle{Retrieval-Augmented Generation (RAG).}  
RAG enhances the response quality of large language models (LLMs) by integrating external knowledge sources during inference, which is particularly beneficial for knowledge-intensive tasks~\cite{huang2025survey}. 
A typical RAG pipeline consists of two stages: \textit{retrieval} and \textit{generation}~\cite{gao2023retrieval}. 
In the retrieval stage, external documents or knowledge bases are first encoded into vector representations and indexed for efficient similarity search~\cite{10.5555/3495724.3496517}. 
Given a query $Q$, the most relevant entities are retrieved and concatenated with $Q$ to form an augmented prompt. 
This enriched prompt is then passed to the LLM, enabling it to generate responses that are more contextually grounded.

\stitle{Chain-of-Thought (CoT).}  
CoT prompting~\cite{wei2022chain} is a widely adopted technique for improving the reasoning capability of LLMs~\cite{chen2025towards}, particularly in complex or multi-step reasoning tasks. 
Instead of directly producing an answer, CoT encourages the model to generate intermediate reasoning steps in a structured format (e.g., \textit{<input, thoughts, output>}). 
These intermediate steps guide the model through reasoning processes such as arithmetic computation, commonsense inference, and symbolic deduction~\cite{wang2022self, lyu2023faithful}. 
Compared to conventional few-shot prompting, CoT enhances both interpretability and performance by decomposing complex reasoning into a sequence of transparent and verifiable sub-tasks.


% \noindent\textbf{LLM-based Agent.}  
% LLM-based agents have demonstrated powerful reasoning and interaction capabilities in real-world settings, spanning from autonomous driving~\cite{mao2023language, yuan2024rag}, knowledge intensive question-answering ~\cite{yue2025survey, yao2022react}, and healthcare~\cite{shi2024ehragent}. These agents, backed by LLM, can take user instructions, gather environmental information, and retrieve knowledge and past experiences from a memory unit to make an informed action plan or tool call.


% \noindent\textbf{Multi Agent System.}  
% Multi-agent systems extend the capabilities of single LLM agents by enabling collaboration and coordination among agents with specialized roles~\cite{talebirad2023multi,zhang2023exploring,park2023generative,li2023camel}. By working together, these agents can tackle tasks that exceed the abilities of any individual agent. In such systems, each agent is assigned distinct capabilities and responsibilities, and they collaborate to achieve shared goals. This interaction—often involving mechanisms like debate, reflection, or role assignment—has been shown to be particularly effective for complex reasoning and creative problem-solving.
% Recent studies have explored multi-agent frameworks in various contexts, including simulated interactive environments~\cite{park2023generative,jinxin2023cgmi}, role-playing~\cite{li2023camel, shao2023character}, and multi-step reasoning~\cite{du2023debate,liang2023encouraging}, demonstrating their strong potential for addressing real-world challenges.

% 这一段是否需要加上?
% RAG-based CoT~\cite{jin2024graph} combines CoT reasoning with retrieval from external knowledge sources. During the reasoning process, the model may fetch relevant information from a RAG system to support multi-hop or hierarchical reasoning. This integration can be guided by predefined rules~\cite{gutierrez2024hipporag, hu2024grag, li2023graph} or dynamically managed by LLM-based agents that decide when and what to retrieve. By blending internal reasoning with external context, RAG-based CoT significantly improves the model's ability to tackle complex, knowledge-intensive tasks with deeper contextual understanding.

% \subsection{RAG}
% While completing tasks through prompting, large language models (LLMs) can often produce incorrect or irrelevant outputs, known as hallucinations\cite{huang2025survey}. Retrieval-Augmented Generation (RAG) addresses this by combining external knowledge with LLMs to improve accuracy and reliability, especially by incorporating domain-specific information. A typical RAG workflow includes two main steps\cite{gao2023retrieval}: (1) \textit{retrieval} — given a user question $Q$, the system searches an indexed corpus (split into smaller chunks) to find the most relevant pieces of information; and (2) \textit{generation} — the LLM then uses the retrieved chunks along with $Q$ as a prompt to generate a response. The external knowledge used in RAG can come in various formats, such as text\cite{izacard2021unsupervised}, images\cite{saito2023pic2word}, videos\cite{arefeen2024vita}, audio\cite{yuan2024retrieval} or code\cite{li2023skcoder}. It also varies in structure, ranging from unstructured plain text\cite{zhu2024realm} to semi-structured HTML\cite{yang2024method} or fully structured graph data\cite{peng2024graph}.

% Retrieval-Augmented Generation (RAG) is a widely adopted technique for improving the generation quality of large language models (LLMs) by incorporating external knowledge into the inference process~\cite{huang2025survey}. This enhancement is especially useful for domain-specific queries, where external context increases accuracy and reliability.

% A typical RAG workflow consists of two stages: \textit{retrieval} and \textit{generation}~\cite{gao2023retrieval}. In the retrieval stage, external knowledge sources—such as documents, images, or code—are embedded into high-dimensional vectors using pre-trained encoders and indexed in a vector database optimized for similarity search. Given a user query $ Q $, the system retrieves the most relevant chunks via semantic search. These retrieved items are then combined with $ Q $ to form an augmented prompt, which is passed to the LLM. By leveraging both the query and the contextual knowledge, the LLM produces responses that are more accurate and context-aware.

% RAG can incorporate various types of external knowledge, including text~\cite{izacard2021unsupervised}, images~\cite{saito2023pic2word}, videos~\cite{arefeen2024vita}, audio~\cite{yuan2024retrieval}, and code~\cite{li2023skcoder}. The data structure of the knowledge also varies—from unstructured plain text~\cite{zhu2024realm}, to semi-structured formats like HTML~\cite{yang2024method}, to fully structured graphs~\cite{peng2024graph}. Graph RAG is a specialized form of RAG where knowledge is organized as structured graphs, capturing not only entities but also their relationships. For example, bibliographic graphs~\cite{wang2020microsoft} represent academic papers and their citation links, revealing collaborations and influence between authors.


% \subsection{Chain-of-Thought}
% Chain-of-Thought (CoT)~\cite{wei2022chain} prompting is a widely used technique that enhances the reasoning capabilities of large language models(LLMs)~\cite{chen2025towards}, particularly for long-text and logic-intensive tasks. Instead of producing a direct answer, CoT encourages the model to generate intermediate reasoning steps in a structured format, e.g., typically \textit{<input, thoughts, output>}. These natural language "thoughts" guide the model through complex tasks such as arithmetic, commonsense reasoning, and symbolic logic~\cite{wang2022self, lyu2023faithful}. Compared to standard few-shot learning, CoT provides greater transparency and improves performance by decomposing problems into smaller, more interpretable steps.

% RAG-based CoT~\cite{jin2024graph} combines CoT reasoning with retrieval from external knowledge sources. During the reasoning process, the model may fetch relevant information from a RAG system to support multi-hop or hierarchical reasoning. This integration can be guided by predefined rules~\cite{gutierrez2024hipporag, hu2024grag, li2023graph} or dynamically managed by LLM-based agents that decide when and what to retrieve. By blending internal reasoning with external context, RAG-based CoT significantly improves the model's ability to tackle complex, knowledge-intensive tasks with deeper contextual understanding.

% \subsection{Problem definition}
% \begin{definition}[Graph--CoT]\label{def:graph-cot}
% Let $G=(V,E)$ be a finite directed graph with $E\subseteq V\times V$.
% Let $\Sigma$ be a token alphabet, $\mathcal{R}$ a finite set of callable operations on $G$ and $\mathcal{Y}$ an answer space.

% Given a natural-language question $q$, the \emph{Graph--CoT} problem is to choose a policy $\pi$ that interleaves chain of thought over LLM-based agent and retrieval calls over $G$ to produce a final answer $y\in\mathcal{Y}$.
% Formally, $\pi$ induces a finite trajectory
% \[
% \tau=(s_0,a_0,s_1,a_1,\ldots,s_T),
% \]

% where each state $s_t=(q,\mathcal{K}_t)$ consists of the question $q$ and an ordered reasoning history $\mathcal{K}_t$, and each action $a_t$ is either:
% \begin{itemize}
% \item Reason: $a_t=\textsc{Reason}(z)$ with $z\in\Sigma^*$, a textual thought produced by an LLM-based agent conditioned on $s_t$.
% \item Retrieve: $a_t=\textsc{Retrieve}(r)$ with $r\in\mathcal{R}$, a retrieval operation to be executed on $G$.
% % its observation is $o=\mathrm{Ret}(r,G)\in\mathcal{O}$.
% \item Finish: $a_t=\textsc{Finish}(y)$ with $y\in\mathcal{Y}$, which outputs the final answer and terminates.
% \end{itemize}

% Let $\log(\cdot)$ denote the logging operator that records the logs of an action, The state update is defined by appending logs: $$\mathcal{K}_{t+1}=\mathcal{K}_t \,\|\, \log(a_t),\qquad s_{t+1}=(q,\mathcal{K}_{t+1})$$

% \noindent\textbf{Objective.} The objective of Graph-CoT problem is to maximize task utility (e.g., answer correctness) while keeping the overall token usage and end-to-end latency as low as possible, without distinguishing between model generation and retrieval costs.
% \end{definition}
% \noindent\textbf{Related Work.}



\subsection{LLM Inference and Prefix Caching}\label{sec:kv_cache_background}

During LLM inference, computation proceeds in two stages: the \textbf{Prefill} stage and the \textbf{Decoding} stage. 
In the prefill stage, the entire input sequence $X = [x_1, x_2, \ldots, x_s]$ is processed to initialize the key–value (KV) cache. 
At each transformer layer, token embeddings are projected into query ($Q$), key ($K$), and value ($V$) representations. 
The $K$ and $V$ tensors are stored in the KV cache for efficient reuse during subsequent decoding. 
Once the first output token $x_{s+1}$ is produced, the model transitions into the autoregressive decoding phase, generating tokens sequentially based on previously generated ones while incrementally updating the KV cache.

Prefix caching is a widely adopted optimization technique for reducing redundant computation in LLM inference. 
The key idea is to cache the KV-cache blocks corresponding to previously processed input prefixes and reuse these cached blocks when new requests share identical prefixes, therefore reducing latency and GPU usage while preserving outputs. 
When GPU memory is limited, cached blocks are evicted according to predefined policies, with Least Recently Used (LRU) being the most common strategy.



% Although decoding is often considered the primary contributor to inference latency, the prefill stage can become a major bottleneck when handling long prompts which are common in RAG workloads. For instance, encoding a 4,000-token prompt using a LLaMA-34B model on an A40 GPU can take over 3 seconds, significantly degrading user experience and limiting system throughput. Recent studies have shown that optimizing or bypassing the prefill stage can nearly double inference throughput\cite{kamath2025pod, shi2025nexus}.


% Although RAG and CoT can improve generation quality, LLM inference performance, including latency and throughput, remains a critical challenge in real-world real-time applications such as OpenAI's O1~\cite{openai2024openaio1card}. To address this, modern LLMs typically adopt a disaggregated key-value (KV) cache framework~\cite{qin2025mooncake}, which separates inference into two sequential stages for efficient KV cache management: the \textit{prefill stage} and the \textit{decoding stage}. These stages are executed in order to process input prompts and generate tokens in an autoregressive manner.

% During LLM inference, the process begins with a \emph{prefill} stage, where the entire input sequence $ X = [x_1, x_2, \ldots, x_s] $ is processed to initialize the key-value (KV) cache. At each transformer layer, the input embeddings are projected into query (Q), key (K), and value (V) vectors, with the K and V outputs cached for efficient reuse during the subsequent \emph{decoding} phase. Once the initial output token $ x_{s+1} $ is generated, the model transitions into autoregressive decoding, producing one token at a time conditioned on prior tokens, while incrementally updating the KV cache. Although decoding is often considered the primary contributor to inference latency, the prefill stage can become a major bottleneck when handling long prompts which are common in RAG workloads. For instance, encoding a 4,000-token prompt using a LLaMA-34B model on an A40 GPU can take over 3 seconds, significantly degrading user experience and limiting system throughput. Recent studies have shown that optimizing or bypassing the prefill stage can nearly double inference throughput\cite{kamath2025pod, shi2025nexus}.

% \textbf{KV Cache Management.}  
% The key-value (KV) cache is a critical optimization in LLM inference that avoids recomputing attention for previously processed tokens. In the prefill stage, a prefix KV cache sharing mechanism allows multiple requests with a common prompt prefix to reuse cached values, reducing redundant computation and memory usage. During the decoding stage, the KV cache stores key and value tensors at each step, enabling efficient token-by-token generation. Without KV cache reuse, the model would need to recompute attention over the entire sequence at every step, leading to substantial computational overhead.

% In addition to reuse, KV cache scheduling is essential for managing memory consumption. KV caching introduces significant memory overhead, e.g., GPT-3 produces approximately 4.5MB of KV cache per token while GPU high-bandwidth memory (HBM) is limited and cannot hold the full cache for all concurrent sequences. To address this, existing systems like vLLM~\cite{kwon2023efficient} typically apply LRU (Least Recently Used) strategies to evict less active cache entries and preserve the most recently accessed ones.


